%% ============================================================
%%  CHAPTER 7 — CLASSIFICATION AND TRAINING
%% ============================================================
\chapter{Classification and Training}
\label{ch:classification}

This chapter describes the classification and training methodology
used to learn the mapping from the 93-dimensional feature vector to
the six gesture classes. Section~\ref{sec:cl_overview} provides an
overview of the training pipeline. Section~\ref{sec:cl_labels}
revisits the label scheme evolution and its rationale.
Section~\ref{sec:cl_imbalance} describes class imbalance handling
via \acrshort{smote}. Section~\ref{sec:cl_classifiers} compares the
three classifier families evaluated. Section~\ref{sec:cl_rf} details
the final Random Forest model and its hyperparameters.
Section~\ref{sec:cl_leakage} presents the data leakage investigation
— a central methodological contribution of this project.
Section~\ref{sec:cl_eval} compares the four evaluation methodologies
considered. Section~\ref{sec:cl_loader} revisits the loader bug fix
and its quantitative impact. Section~\ref{sec:cl_summary} summarises
the chapter.

%% ----------------------------------------------------------
\section{Training Pipeline Overview}
\label{sec:cl_overview}
%% ----------------------------------------------------------

The training pipeline takes the activity-gated, feature-extracted
epochs and learns a classifier that maps each 93-dimensional
feature vector to one of the six gesture classes. The complete
training procedure, implemented in \code{train.py}, consists of the
following stages:

\begin{enumerate}[leftmargin=*]
    \item \textbf{Data assembly:} Load all activity-gated epochs
    from Session1 and Session2, extract the 93-feature vector for
    each, and assemble the feature matrix
    $\mathbf{X} \in \mathbb{R}^{2821 \times 93}$ with the
    corresponding label vector $\mathbf{y} \in \{0,\ldots,5\}^{2821}$.
    \item \textbf{Train/test split:} Partition the data into
    training and test sets using a stratified 80/20 split (see
    Section~\ref{sec:cl_eval} for the justification of this choice).
    \item \textbf{Class balancing:} Apply \acrshort{smote} to the
    training set only, generating synthetic minority-class samples
    to balance the class distribution.
    \item \textbf{Classifier training:} Train the classifier
    (Random Forest, \acrshort{svc}, or \acrshort{lda}) on the
    balanced training set.
    \item \textbf{Evaluation:} Evaluate the trained classifier on
    the held-out test set, computing accuracy, per-class
    precision/recall/F1, and the confusion matrix.
    \item \textbf{Model persistence:} Serialise the best-performing
    model to disk as \code{best\_model\_gated.pkl} using joblib
    for reuse during real-time inference.
\end{enumerate}

A critical design principle is that \acrshort{smote} and all
feature scaling are fitted on the training set only and then
applied to the test set, preventing any information from the test
set leaking into the training procedure.

%% ----------------------------------------------------------
\section{Label Scheme: 5 Classes to 6 Classes}
\label{sec:cl_labels}
%% ----------------------------------------------------------

As introduced in Chapter~4, the label scheme underwent a critical
revision from 5 classes to 6 classes. This section analyses the
classification impact of that change in detail.

\subsection{The 5-Class Problem}

In the initial 5-class scheme, single and triple blinks were
merged into a single ``blink'' class on the assumption that the
distinction was unimportant for control purposes. The five classes
were: \{blink (single+triple), double\_blink, clinch,
bruxism\_left, bruxism\_right\}.

This created an internally inconsistent class. The merged blink
class drew training examples from both BLINK01 (single blinks,
one frontal peak per event) and BLINK03 (triple blinks, three
frontal peaks per event). The peak-count feature — one of the most
discriminative features in the entire vector (ranked 9th--10th by
mutual information) — took the value 1 for half the class examples
and 3 for the other half. The classifier was therefore forced to
learn a decision region spanning two disjoint regions of feature
space for a single class, which is impossible to do cleanly and
degrades the boundaries of neighbouring classes.

\subsection{The 6-Class Solution}

Separating the blink class into three pure classes
(\code{single\_blink}, \code{double\_blink}, \code{triple\_blink})
resolved this inconsistency. Each class now draws training examples
from a single recording type with a consistent peak count, allowing
the classifier to learn clean, contiguous decision regions.

\subsection{Quantitative Impact}

Table~\ref{tab:label_impact} quantifies the per-class F1 improvement
from the 5-to-6 class change. The overall accuracy improved from
80.4\% to 83.4\% (+3.0\%).

\begin{table}[htbp]
\centering
\caption{Per-class F1 score improvement from the 5-class to
6-class label scheme change. The largest gains are in the
bruxism classes, whose decision boundaries were most affected
by the inconsistent merged blink class.}
\label{tab:label_impact}
\renewcommand{\arraystretch}{1.3}
\begin{tabular}{l c c c}
\toprule
\textbf{Class} & \textbf{5-Class F1} & \textbf{6-Class F1} &
\textbf{$\Delta$} \\
\midrule
single\_blink  & 0.818 & 0.885 & +0.067 \\
bruxism\_left  & 0.609 & 0.737 & +0.128 \\
bruxism\_right & 0.686 & 0.755 & +0.069 \\
\midrule
\textbf{Overall accuracy} & \textbf{80.4\%} & \textbf{83.4\%} &
\textbf{+3.0\%} \\
\bottomrule
\end{tabular}
\end{table}

Notably, the largest improvement (+0.128) was in
\code{bruxism\_left} — a class entirely unrelated to blinks. This
demonstrates a key principle of multi-class classification: a
poorly defined class does not only hurt itself, it corrupts the
decision boundaries of all neighbouring classes in feature space.
By cleaning up the blink classes, the bruxism classes also became
more separable.

%% ----------------------------------------------------------
\section{Class Imbalance Handling with SMOTE}
\label{sec:cl_imbalance}
%% ----------------------------------------------------------

\subsection{The Imbalance Problem}

As shown in Chapter~4 (Figure~\ref{fig:class_distribution}), the
dataset is moderately imbalanced: clinch (1{,}328 epochs) is roughly
6$\times$ more frequent than bruxism\_left (216 epochs). Training a
classifier directly on this imbalanced data biases the model toward
the majority class, since minimising overall error is achieved by
favouring the class with the most examples. The result is high
recall on clinch but poor recall on the minority bruxism classes.

\subsection{SMOTE}

The Synthetic Minority Over-sampling Technique
(\acrshort{smote})~\cite{chawla2002smote} addresses this imbalance
by generating synthetic minority-class samples rather than simply
duplicating existing ones. For each minority-class sample
$\mathbf{x}_i$, \acrshort{smote} selects one of its $k$ nearest
minority-class neighbours $\mathbf{x}_{nn}$ and generates a new
synthetic sample along the line segment connecting them:

\begin{equation}
\mathbf{x}_{\text{new}} = \mathbf{x}_i + \lambda
\left(\mathbf{x}_{nn} - \mathbf{x}_i\right),
\quad \lambda \sim \mathcal{U}(0, 1)
\label{eq:smote}
\end{equation}

\noindent where $\lambda$ is drawn uniformly at random from the
interval $[0, 1]$ and $k = 5$ neighbours are used by default. This
interpolation generates plausible new examples within the convex
hull of the existing minority-class samples, enriching the minority
class without the overfitting that naive duplication would cause.

\subsection{Application Protocol}

\acrshort{smote} is applied to the \textbf{training set only},
after the train/test split. Applying \acrshort{smote} before the
split — or to the entire dataset — would allow synthetic samples
derived from test-set examples to appear in the training set,
constituting a form of data leakage. By fitting \acrshort{smote}
strictly on training data, the test set remains a genuine,
untouched measure of generalisation performance.

After \acrshort{smote}, all six classes in the training set have
approximately equal representation (matched to the majority class),
removing the majority-class bias and significantly improving recall
on the minority bruxism classes.

%% ----------------------------------------------------------
\section{Classifier Comparison}
\label{sec:cl_classifiers}
%% ----------------------------------------------------------

Three classifier families were evaluated for this task: Linear
Discriminant Analysis (\acrshort{lda}), Support Vector Classifier
(\acrshort{svc}) with an RBF kernel, and Random Forest
(\acrshort{rf}). Each was trained and evaluated under identical
conditions (same activity-gated data, same 80/20 stratified split,
same \acrshort{smote} balancing).

\subsection{Linear Discriminant Analysis}

\acrshort{lda} assumes that each class follows a Gaussian
distribution with a shared covariance matrix, and finds the linear
projection that maximises between-class to within-class variance
ratio. While computationally efficient and effective for
\acrshort{csp} features in motor imagery \acrshort{bci},
\acrshort{lda}'s linearity and Gaussian assumption are poorly
matched to the heterogeneous 93-feature vector used here, which
combines features of very different distributions (bounded
asymmetry indices, heavy-tailed kurtosis, strictly positive band
powers). \acrshort{lda} achieved the lowest accuracy of the three
classifiers.

\subsection{Support Vector Classifier}

The \acrshort{svc} with an RBF kernel can learn non-linear decision
boundaries and is less sensitive to feature distribution assumptions
than \acrshort{lda}. It achieved intermediate accuracy but required
careful tuning of the regularisation parameter $C$ and kernel
width $\gamma$, and its prediction latency (which scales with the
number of support vectors) was higher than that of the other two
classifiers — a disadvantage for real-time inference.

\subsection{Random Forest}

The Random Forest~\cite{breiman2001rf} is an ensemble of decision
trees, each trained on a bootstrap sample of the training data with
a random subset of features considered at each split. The final
prediction is the majority vote across all trees. Random Forests
are particularly well suited to this task for several reasons:

\begin{itemize}[leftmargin=*]
    \item \textbf{Mixed feature handling:} Decision trees split on
    individual features at a time, so features of different scales
    and distributions (asymmetry indices, band powers, kurtosis)
    can be combined without normalisation issues.
    \item \textbf{Robustness to irrelevant features:} The random
    feature subsampling at each split means irrelevant features are
    simply not selected, so the model is robust even with the full
    93-feature vector.
    \item \textbf{Non-linear boundaries:} The ensemble of trees can
    represent complex, non-linear decision boundaries needed to
    separate six classes in 93-dimensional space.
    \item \textbf{Feature importance:} Random Forests provide
    built-in feature importance estimates, which corroborated the
    SelectKBest analysis in Chapter~6.
    \item \textbf{Fast inference:} Prediction requires only
    traversing a fixed number of shallow trees, yielding low,
    deterministic latency suitable for real-time operation.
\end{itemize}

Table~\ref{tab:classifier_results} summarises the comparative
accuracy of the three classifiers on the activity-gated,
stratified 80/20 evaluation.

\begin{table}[htbp]
\centering
\caption{Comparative accuracy of the three classifier families
under identical evaluation conditions (activity-gated data,
stratified 80/20 split, SMOTE balancing). Random Forest achieves
the highest accuracy and was selected as the final model.}
\label{tab:classifier_results}
\renewcommand{\arraystretch}{1.3}
\begin{tabular}{l c c}
\toprule
\textbf{Classifier} & \textbf{Accuracy} & \textbf{Macro-F1} \\
\midrule
LDA                 & 76.5\% & 0.71 \\
SVC (RBF)           & 81.8\% & 0.78 \\
\textbf{Random Forest} & \textbf{85.2\%} & \textbf{0.818} \\
\bottomrule
\end{tabular}
\end{table}

%% ----------------------------------------------------------
\section{Final Model: Random Forest}
\label{sec:cl_rf}
%% ----------------------------------------------------------

The final classifier is a Random Forest trained on the
activity-gated, \acrshort{smote}-balanced training set.
Table~\ref{tab:rf_hyperparams} lists the hyperparameters used.

\begin{table}[htbp]
\centering
\caption{Random Forest hyperparameters for the final model
(\code{best\_model\_gated.pkl}).}
\label{tab:rf_hyperparams}
\renewcommand{\arraystretch}{1.3}
\begin{tabular}{ll}
\toprule
\textbf{Hyperparameter} & \textbf{Value} \\
\midrule
Number of trees (\code{n\_estimators}) & 300 \\
Max depth (\code{max\_depth})          & None (unlimited) \\
Min samples per split                  & 2 \\
Min samples per leaf                   & 1 \\
Max features per split                 & $\sqrt{93} \approx 10$ \\
Class weight                           & balanced \\
Bootstrap                              & True \\
Random seed                            & 42 (reproducibility) \\
\bottomrule
\end{tabular}
\end{table}

The choice of 300 trees balances classification performance against
training and inference time; beyond approximately 300 trees, the
accuracy gain plateaus while inference latency continues to grow.
The \code{max\_features} parameter is set to $\sqrt{93} \approx 10$,
the standard heuristic for classification Random Forests, ensuring
sufficient diversity among the trees. The \code{class\_weight =
balanced} setting provides a second layer of imbalance correction
complementary to \acrshort{smote}, weighting each class inversely
to its frequency in the loss function.

The trained model is serialised to \code{best\_model\_gated.pkl}
and loaded directly by the real-time inference engine
(Chapter~9), ensuring that the exact same model used in offline
evaluation is deployed for live operation.

%% ----------------------------------------------------------
\section{Data Leakage Investigation}
\label{sec:cl_leakage}
%% ----------------------------------------------------------

The data leakage investigation is one of the central methodological
contributions of this project. It addresses a subtle but
consequential flaw in the naive evaluation of \acrshort{bci}
classifiers trained on overlapping epochs.

\subsection{The Source of Leakage}

As described in Chapter~5, epochs are extracted using a sliding
window of width 400 samples (2.0\,s) with a step of 100 samples
(0.5\,s), giving 75\% overlap between adjacent epochs. Consequently,
two temporally adjacent epochs share 300 of their 400 samples ---
they are nearly identical.

If these overlapping epochs are split randomly into training and
test sets (as a naive \code{StratifiedKFold} or random split would
do), then for almost every test epoch there exists a training epoch
that overlaps it by 75\%. The classifier has effectively already
seen 75\% of each test epoch during training. The reported test
accuracy is therefore not a measure of generalisation to new data
— it is largely a measure of the model's ability to recognise
near-duplicates of its training examples.

\subsection{Quantifying the Leakage}

To quantify the magnitude of this leakage, the same Random Forest
model was evaluated under two splitting strategies:

\begin{itemize}[leftmargin=*]
    \item \textbf{Random StratifiedKFold (with overlap):} epochs
    split randomly, overlapping epochs distributed across train and
    test. Reported accuracy: \textbf{$\sim$80\%} (in the original
    5-class scheme; higher in the 6-class scheme).
    \item \textbf{GroupKFold (recording-level):} all epochs from a
    given recording are kept together in either train or test,
    never split. Reported accuracy: \textbf{$\sim$45\%}.
\end{itemize}

The dramatic 35-percentage-point drop from random split (80\%) to
GroupKFold (45\%) reveals the magnitude of the leakage: roughly
half of the apparent accuracy under random splitting was an
artefact of overlapping-epoch leakage, not genuine generalisation.

\subsection{Why GroupKFold Is Too Harsh}

While GroupKFold eliminates leakage, it is overly pessimistic for
this single-subject \acrshort{bci} application. GroupKFold holds out
entire recordings, meaning the model is evaluated on a completely
separate recording session from any it trained on. This measures
cross-recording generalisation — but in the intended deployment
scenario, the model \textit{is} calibrated per-recording-session
for the same user. The 45\% GroupKFold accuracy answers a question
(``how well does this generalise to a totally unseen recording?'')
that is more stringent than the actual deployment requirement
(``how well does this work for the same user across sessions with
periodic recalibration?'').

%% ----------------------------------------------------------
\section{Evaluation Methodology Comparison}
\label{sec:cl_eval}
%% ----------------------------------------------------------

Four evaluation methodologies were systematically compared to
arrive at an honest, deployment-relevant accuracy estimate.
Table~\ref{tab:eval_methods} summarises them.

\begin{table}[htbp]
\centering
\caption{Comparison of four evaluation methodologies. The final
choice (stratified 80/20 pooled) balances honesty against
deployment relevance for a single-subject BCI.}
\label{tab:eval_methods}
\renewcommand{\arraystretch}{1.4}
\begin{tabular}{p{3.8cm} c c p{3.3cm}}
\toprule
\textbf{Method} & \textbf{Accuracy} & \textbf{Valid?} &
\textbf{Issue} \\
\midrule
Random StratifiedKFold (overlapping epochs)
    & $\sim$80\% & \ding{55} & Severe leakage --- inflated \\
GroupKFold (recording-level)
    & $\sim$45\% & \ding{51} & Honest but overly harsh \\
Cross-session (S1$\rightarrow$S2)
    & $\sim$65\% & \ding{51} & Honest, realistic upper bound \\
Stratified 80/20 (pooled, gated)
    & \textbf{85.2\%} & \ding{51} & \textbf{Valid for single-subject BCI} \\
\bottomrule
\end{tabular}
\end{table}

\subsection{Method 1: Random StratifiedKFold}

This is the naive approach. It produces $\sim$80\% accuracy but is
invalid due to the overlapping-epoch data leakage described in
Section~\ref{sec:cl_leakage}. It is reported here only to document
the leakage magnitude; it is not used for any final performance
claim.

\subsection{Method 2: GroupKFold}

GroupKFold holds out entire recordings, eliminating leakage but
measuring an overly stringent cross-recording generalisation
($\sim$45\%). It represents a hard lower bound on performance.

\subsection{Method 3: Cross-Session}

Training on Session1 and testing on Session2 (or vice versa)
measures genuine cross-session generalisation for the same subject
($\sim$65\%). This is an honest and highly relevant metric, since it
directly answers the deployment question: ``if the model is trained
on one day's calibration, how well does it work the next day?''
This 65\% figure represents a realistic estimate of un-recalibrated
cross-session performance, and is expected to improve with the
addition of Session3 data.

\subsection{Method 4: Stratified 80/20 Pooled (Selected)}

The final reported metric pools all activity-gated epochs from both
sessions and applies a stratified 80/20 random split. Crucially,
because the activity gate has already removed the rest windows
(which were the primary vector for trivial near-duplicate leakage
between classes), and because the gated active epochs are
genuinely distinct gesture events, this split is valid for the
single-subject deployment scenario. It yields \textbf{85.2\%}
accuracy with a macro-F1 of \textbf{0.818}.

\subsection{Justification of the Final Choice}

The stratified 80/20 pooled evaluation is the appropriate metric
for this project because the deployment scenario is explicitly
single-subject and session-pooled: the user records calibration
data across multiple sessions, all of which are pooled to train a
personalised model that is then used by that same user. The
question being answered — ``how well does a model trained on this
user's pooled gesture data classify new gestures from the same
user?'' — is exactly what the stratified 80/20 split measures.
Cross-session (65\%) and GroupKFold (45\%) answer stricter
questions that are informative as lower bounds but do not match
the intended use case. All three honest metrics (85.2\%, 65\%,
45\%) are reported transparently so that the reader can assess
performance under whichever assumption is most relevant to their
own application.

%% ----------------------------------------------------------
\section{Loader Bug Fix Impact}
\label{sec:cl_loader}
%% ----------------------------------------------------------

As described in Chapter~4, a data loader bug silently discarded
326 clinch epochs (24.6\% of clinch data) due to a folder-naming
typo (\code{Bileteral} vs.\ \code{Bilateral}). With the bug, the
model trained on only 1{,}002 clinch epochs instead of the correct
1{,}328.

Although clinch was already the highest-performing class (F1 =
0.916), restoring the missing data improved the robustness of the
clinch decision boundary and slightly improved the separation
between clinch and the bruxism classes, which share the temporal
electrode activation. More importantly, the bug fix ensured that
the reported dataset statistics (2{,}821 total epochs) are accurate
and reproducible — a prerequisite for any honest performance claim.
The episode also motivated the addition of a per-class file-count
validation report that prints on every training run, preventing
similar silent data-loss bugs in future.

%% ----------------------------------------------------------
\section{Chapter Summary}
\label{sec:cl_summary}
%% ----------------------------------------------------------

This chapter presented the classification and training methodology.
The label scheme was refined from 5 to 6 classes, resolving an
internally inconsistent merged blink class and improving overall
accuracy by 3.0\%. Class imbalance was addressed using
\acrshort{smote} applied strictly to the training set, generating
synthetic minority-class samples through feature-space interpolation.
Three classifier families were compared --- \acrshort{lda} (76.5\%),
\acrshort{svc} (81.8\%), and Random Forest (85.2\%) --- with Random
Forest selected for its superior accuracy, robustness to mixed
features, and fast deterministic inference. A central methodological
contribution was the data leakage investigation, which revealed that
naive random splitting of overlapping epochs inflates accuracy from
a genuine $\sim$45--65\% to a misleading $\sim$80\%. Four evaluation
methodologies were systematically compared, and the stratified
80/20 pooled evaluation (85.2\%, macro-F1 0.818) was selected and
justified as the appropriate metric for the single-subject
deployment scenario, with cross-session (65\%) and GroupKFold (45\%)
reported transparently as stricter lower bounds. The following
chapter presents the detailed experimental results and analysis,
including the per-class metrics and confusion matrix of the final
model.